\documentclass[11pt,a4paper,openany,dvipsnames]{book}
\pagestyle{plain}

\usepackage{../core-rulebook/config/packages}
\usepackage{../core-rulebook/config/news}

% Title definition
\def\subtitle{Arsenal}

\begin{document}
{
	\heading
	\pagenumbering{gobble}
\newgeometry{margin=0.3cm}

\newlength{\neontitlewidth}
\settowidth{\neontitlewidth}{NEON TTRPG}
\newlength{\neonsubtitlewidth}
\settowidth{\neonsubtitlewidth}{\subtitle}
\def\coverborderthickness{1mm}
\def\coverbordercolor{ProcessBlue}

\begin{center}
\IfFileExists{\coverimgpath}{\ThisCenterWallPaper{1}{\coverimgpath}}{}
\begin{tikzpicture}
	
	%conf
	\coordinate (toptextanchor)
		at ( $ (current page.north) - (7.3,1) $ )
	;
	\coordinate (bottomtextanchor)
		at ( $ (current page.south east) - (3.3,-1) $ )
	;
	
	%TITLE
	\node [white,
		scale=5,
		anchor=west,
		inner sep=0pt,
	]
		at ( $ (toptextanchor) - (0,0.3) $ )
		{NEON TTRPG}
	;
	%SUBTITLE
	\node [white,
		scale=4,
		anchor=east,
		inner sep=0pt,
	]
		at ( $ (bottomtextanchor) + (0,0.3) $ )
		{\subtitle}
	;
	
	%LEFT
	\draw [\coverbordercolor,
		line width=\coverborderthickness,
	]
	(toptextanchor)
	-- ( $ (current page.north west) + (3,-1) $)
	-- ++(-2,-2)
	-- ( $ (current page.south west) + (1,3) $ )
	-- ++(2,-2)
	-- ( $ (bottomtextanchor) - (4\neonsubtitlewidth,0) $)
	;
	
	%RIGHT
	\draw [\coverbordercolor,
		line width=\coverborderthickness,
	]
	( $ (toptextanchor) + (5\neontitlewidth,0) $)
	-- ( $ (current page.north east) - (3,1) $)
	-- ++(2,-2)
	-- ( $ (current page.south east) - (1,-3) $ )
	-- ++(-2,-2)
	-- (bottomtextanchor)
	;
	
\end{tikzpicture}

\end{center}\restoregeometry

}
{\hypersetup{hidelinks} \tableofcontents}


%% RANGED
\chapter{Ranged}
\section{Weapons}
    \luaimport{equipment/rangedweapons.json}{ranged-arsenal.tpl}{equipment/ranged}

\section{Attachments}
\section{Ammunition}
	\intablefalse
	\luaimport{equipment/ammo.json}{ammo-arsenal.tpl}{equipment/ammo}
	\end{tabular}

%% MELEE
\chapter{Melee}
\section{Assembly}
Melee weapons are not set in stone.
Instead they are created from various components, allowing for ideal customization.\\
First, any melee weapon needs a grip.
Swords, knives and whips have \emph{hilts}, while hammers, axes and spears have \emph{shafts}.\\
The weapon point (the "business end") is the second component.
Swords, knives and spears have \emph{blades};
	axes and hammers have \emph{heads} and \emph{whips} each have their own category.\\
Lastly every melee weapon may also have a \emph{core} component.
These are never required, they but add unique features to the weapon.
\section{Base Damage}
Melee weapons deal damage based on their size.
\begin{sitemize}
	\item \emph{Unarmed fists} deal \fbox{1D5 + the user's Melee Damage Bonus}, and are \emph{Improvised}.
	\item \emph{One-handed melee weapons} deal \fbox{1D10 + the user's Melee Damage Bonus}.
	\item \emph{Two-handed weapons} deal \fbox{2D10 + the user's Melee Damage Bonus}.
	\item \emph{Oversized melee weapons}, that cannot be used by humans anymore, deal \fbox{4D10}.
		Their damage may be increased by the contraption wielding it.
\end{sitemize}
\pagebreak%quick fix
\section{Components}
    \subsection{Hilts}
    \vspace{8mm}
    \begin{multicols}{2}
	    \luaimport{equipment/meleecomponents-hilt.csvin}{meleecomponent.tpl}{equipment/meleecomponent-hilt}
    \end{multicols}

    \subsection{Shafts}
    \vspace{8mm}
    \begin{multicols}{2}
    	\luaimport{equipment/meleecomponents-shaft.csvin}{meleecomponent.tpl}{equipment/meleecomponent-shaft}
    \end{multicols}

    \subsection{Blades}
    \vspace{8mm}
    \begin{multicols}{2}
    	\luaimport{equipment/meleecomponents-blade.csvin}{meleecomponent.tpl}{equipment/meleecomponent-blade}
    \end{multicols}

    \subsection{Whips}
    \vspace{8mm}
    \begin{multicols}{2}
    	\luaimport{equipment/meleecomponents-whip.csvin}{meleecomponent.tpl}{equipment/meleecomponent-whip}
    \end{multicols}

    \subsection{Heads}
    \vspace{8mm}
    \begin{multicols}{2}
    	\luaimport{equipment/meleecomponents-head.csvin}{meleecomponent.tpl}{equipment/meleecomponent-head}
    \end{multicols}

    \subsection{Core}
    \vspace{8mm}
    \begin{multicols}{2}
    	\luaimport{equipment/meleecomponents-core.csvin}{meleecomponent.tpl}{equipment/meleecomponent-core}
    \end{multicols}

%% ARMOR
\chapter{Armor}

\printindex

\IfFileExists{../core-rulebook/config/builddata.texin}{\input{../core-rulebook/config/builddata.texin}}{ }
\end{document}
